\documentclass[11pt]{article}

\usepackage[a4paper,margin=2.5cm]{geometry}
\usepackage{algorithm}
\usepackage{algorithmic}
\usepackage{amsmath}
\usepackage[utf8]{inputenc}
\usepackage[T1]{fontenc}

\renewcommand{\algorithmicrequire}{\textbf{Input:}}
\renewcommand{\algorithmicensure}{\textbf{Output:}}

\title{Risoluzione delle rilocazioni dei simboli esterni in ELF\\
\large (dopo il caricamento di tutte le \texttt{DT\_NEEDED})}
\author{}
\date{}

\begin{document}
\maketitle

\noindent
Notazione: ogni oggetto caricato $obj$ (eseguibile o shared object) espone le
proprie sezioni dinamiche: $obj.\texttt{relaDyn}$ (rilocazioni eager, da
\texttt{.rela.dyn}), $obj.\texttt{relaPlt}$ (rilocazioni delle funzioni, da
\texttt{.rela.plt}), $obj.\texttt{dynsym}$ (tabella dei simboli,
\texttt{.dynsym}), $obj.\texttt{dynstr}$ (tabella delle stringhe,
\texttt{.dynstr}), $obj.\texttt{gnuHash}$ (\texttt{.gnu.hash}),
$obj.\texttt{got}$ / $obj.\texttt{gotPlt}$ (\texttt{.got} e \texttt{.got.plt},
dove viene \emph{scritto} il risultato) e $obj.\texttt{loadBias}$ (base di
caricamento scelta a runtime, es. da ASLR).

\begin{algorithm}[h]
\caption{\textsc{RisolviRilocazioni} --- eseguito da \texttt{ld.so} per ogni
oggetto $obj$ del grafo di dipendenze, dopo che tutte le \texttt{DT\_NEEDED}
sono state mappate in memoria}
\begin{algorithmic}[1]
\REQUIRE $L = [obj_1,\dots,obj_n]$ \COMMENT{tutti gli oggetti caricati:
eseguibile + DT\_NEEDED transitive, gia' mappati in memoria}
\ENSURE per ogni $obj \in L$, $obj.\texttt{got}$ / $obj.\texttt{gotPlt}$
contengono l'indirizzo runtime corretto per ogni simbolo esterno referenziato

\FOR{ogni $obj \in L$}
    \STATE $R \gets obj.\texttt{relaDyn}$ \COMMENT{voci lette da \texttt{.rela.dyn} \emph{di obj}}
    \FOR{ogni rilocazione $r \in R$}
        \IF{$r.\text{tipo} = \texttt{R\_*\_RELATIVE}$}
            \STATE $addr \gets obj.\texttt{loadBias} + r.\text{addend}$ \COMMENT{nessuna ricerca di simbolo: solo aggiustamento per ASLR}
        \ELSE
            \IF{$r.\text{tipo} \in \{\texttt{R\_*\_GLOB\_DAT}, \texttt{R\_*\_JUMP\_SLOT}\}$}
                \STATE $sIdx \gets r.\text{sym\_index}$
                \STATE $sName \gets obj.\texttt{dynstr}[\, obj.\texttt{dynsym}[sIdx].\text{st\_name} \,]$ \COMMENT{nome letto da \texttt{.dynsym}/\texttt{.dynstr} \emph{di obj}, il richiedente}
                \STATE $target \gets \textsc{RisolviSimbolo}(sName, L, obj)$
                \IF{$target = \text{NULL}$}
                    \STATE \textbf{errore}: simbolo indefinito $sName$ \COMMENT{a meno che sia \texttt{STB\_WEAK}: in tal caso $addr \gets 0$}
                \ELSE
                    \STATE $(defObj, defIdx) \gets target$
                    \STATE $addr \gets defObj.\texttt{loadBias} + defObj.\texttt{dynsym}[defIdx].\text{st\_value}$ \COMMENT{indirizzo calcolato nel binario/.so che \emph{definisce} il simbolo}
                \ENDIF
            \ENDIF
        \ENDIF
        \STATE scrivi $addr$ in $obj.\texttt{got}$ (dati) o $obj.\texttt{gotPlt}$ (funzioni) all'offset $r.\text{offset}$ \COMMENT{scrittura sempre nella GOT \emph{di obj}, il richiedente}
    \ENDFOR
\ENDFOR
\end{algorithmic}
\end{algorithm}

\begin{algorithm}[h]
\caption{\textsc{RisolviSimbolo} --- individua \emph{in quale} oggetto del
grafo il simbolo e' definito, seguendo lo scope globale}
\begin{algorithmic}[1]
\REQUIRE nome $name$, lista oggetti $L$, oggetto richiedente $requester$
\ENSURE coppia $(obj, idx)$ del primo simbolo definito trovato, oppure NULL
\FOR{ogni $obj \in L$ \COMMENT{ordine: eseguibile, poi librerie nell'ordine di caricamento}}
    \STATE $idx \gets \textsc{GnuHashLookup}(name,\ obj.\texttt{gnuHash},\ obj.\texttt{dynsym},\ obj.\texttt{dynstr})$
    \IF{$idx \neq \text{NOT\_FOUND}$ \textbf{e} $obj.\texttt{dynsym}[idx].\text{st\_shndx} \neq \texttt{SHN\_UNDEF}$}
        \STATE \textbf{return} $(obj, idx)$ \COMMENT{primo match nello scope: puo' essere anche requester stesso (interposition)}
    \ENDIF
\ENDFOR
\STATE \textbf{return} NULL
\end{algorithmic}
\end{algorithm}

\begin{algorithm}[h]
\caption{\textsc{GnuHashLookup} --- ricerca del simbolo dentro un singolo
oggetto, tramite \texttt{.gnu.hash}}
\begin{algorithmic}[1]
\REQUIRE nome $name$, sezione $hashSec$ (\texttt{.gnu.hash}), $symtab$ (\texttt{.dynsym}), $strtab$ (\texttt{.dynstr}) \emph{tutti dello stesso oggetto candidato}
\ENSURE indice in $symtab$, oppure NOT\_FOUND
\STATE $h \gets \text{GnuHash32}(name)$
\IF{bloom filter di $hashSec$ non contiene $h$}
    \STATE \textbf{return} NOT\_FOUND \COMMENT{scarto immediato, nessun accesso a symtab/strtab}
\ENDIF
\STATE $b \gets h \bmod \text{nbuckets}(hashSec)$
\STATE $i \gets \text{bucket}[b]$
\WHILE{$i \neq 0$}
    \IF{$(\text{chain}[i] \mid 1) = (h \mid 1)$} \COMMENT{confronto hash a 32 bit, economico}
        \IF{$strtab[\, symtab[i].\text{st\_name} \,] = name$} \COMMENT{solo ora si tocca .dynstr}
            \STATE \textbf{return} $i$
        \ENDIF
    \ENDIF
    \IF{$\text{chain}[i]$ ha il bit di stop settato}
        \STATE \textbf{break} \COMMENT{fine della catena del bucket}
    \ENDIF
    \STATE $i \gets i + 1$
\ENDWHILE
\STATE \textbf{return} NOT\_FOUND
\end{algorithmic}
\end{algorithm}

\end{document}
